\chapter{System Overview and Architecture}
\label{ch:system}

\section{4-Layer Bits-to-Atoms Stack}

We propose a 4-layer architecture inspired by principles of hardware-software co-design \cite{Hennessy2019} and abstraction hierarchies in systems design \cite{Simon1996}. This stack separates concerns across computational abstraction levels, enabling reuse and generalization:

\subsection{Layer 1: Physical World}

Robotic agents with sensors, actuators, and communication links operating in dynamic environments. Real-time constraints (10--20 ms decision-action cycles) drive the entire stack's design \cite{Siciliano2016}. This layer generates MRTA problem instances and consumes control commands.

\subsection{Layer 2: Neuromorphic Hardware}

OIM (Oscillator Ising Machines) \cite{Wang2019} for combinatorial optimization, and SNNs (Spiking Neural Networks) \cite{Gerstner2014} for iterative control. Both exploit event-driven, low-power computation principles native to physical systems, enabling sub-milliwatt operation over millisecond timescales. Hardware abstraction shields upper layers from implementation details \cite{Hennessy2019}.

\subsection{Layer 3: Mathematical Encoding}

Compilation layer: MRTA $\to$ MWIS $\to$ QUBO $\to$ Ising (for OIM); robot dynamics $\to$ constrained QP $\to$ iterative algorithm (for SNN). Problem-independent mathematical transformations enable reuse of hardware solvers across applications \cite{Kochenberger2014, Lucas2014}.

\subsection{Layer 4: Problem Formulation}

Application-specific parameters: coalition sizes, task requirements, robot capabilities, MPC horizons. Domain expertise enters at this layer only, preserving generality in lower layers. This separation follows best practices in engineering design \cite{Simon1996}.

\section{OIM vs SNN Use Cases}

The two neuromorphic approaches address complementary problems in multi-agent robotics:

\begin{center}
\begin{tabular}{lll}
\toprule
Property & OIM-MRTA & SNN-MPC \\ \midrule
Problem type & Combinatorial optimization & Continuous control \\
Algorithm & Phase locking \cite{Wang2019} & Iterative gradient projection \cite{He2016} \\
Time horizon & One-shot allocation & Receding horizon ($N=10$ steps) \\
Latency budget & 10--15 ms & 20 ms per cycle \\
Power consumption & $\sim$1 mW (100 oscillators) & $<$1 $\mu$W (sparse spiking) \\
Scalability & $N < 100$ tasks & Fixed QP size (4--10 DOF) \\
Approximation & 92\% of optimal & Converges to optimality \\
\bottomrule
\end{tabular}
\end{center}

\section{Trade-off Space}

Three regimes of hardware-algorithm alignment:

\textbf{Classical Approaches (CPU-based):}
\begin{itemize}
\item Exact MWIS: NP-hard, $O(2^N)$ time \cite{Korsah2012}
\item OSQP MPC: $O(N^3)$ per solve \cite{Stellato2018}, requires milliseconds to tens of milliseconds
\item Energy: 10--100 W continuous, $>$100 mJ per solve
\item Suitable for: offline planning, batch processing, non-real-time robotics
\end{itemize}

\textbf{Neuromorphic Approaches (OIM/SNN):}
\begin{itemize}
\item OIM: Polynomial hardware time (microseconds to milliseconds), $O(N)$ per iteration, approximate solutions \cite{Wang2019}
\item SNN-MPC: Event-driven spiking, convergence in 50--200 ms, sub-milliwatt power \cite{Mangalore2024}
\item Energy: milliwatts to microwatts (1000--100,000$\times$ improvement)
\item Suitable for: real-time robotics, embedded autonomous systems, energy-constrained platforms
\end{itemize}

OIM and SNN excel precisely when latency and energy are critical constraints---the regime of fast-moving robotic teams in the field. This is the domain where neuromorphic hardware offers transformative advantages, justifying the added complexity of specialized hardware.

